%%
%% winter-lecture28.tex
%% 
%% Made by Alex Nelson <pqnelson@gmail.com>
%% Login   <alex@lisp>
%% 
%% Started on  2026-03-10T09:09:01-0700
%% Last update 2026-03-10T09:09:01-0700
%% 

\lecture[Artin--Schreier Theorem]{}

Today and next lecture are student presentations on topics.

\begin{node}
When is it possible to have an algebraic closure of a field which is a
finite extension? The Artin--Schreier theorem says the only way to do
this is by adjoining a quadratic extension to a field. For example
$\CC=\RR(\I)$.
\end{node}

\begin{theorem}[Artin--Schreier]
Let $L/F$ be a field extension.
If $L$ is an algebraically closed field with $1<[L:F]<\infty$, then
$L=F(\I)$ where $\I=\sqrt{-1}$ and $\Char(F)=0$.
\end{theorem}

Before proving this, let us see a quick consequence.

\begin{corollary}
If $\RR/K$ is finite, then $K=\RR$.
\end{corollary}

\begin{proof}
We see $\CC/K$ is finite by Artin--Schreier, so $\CC=K(\I)$, so
\begin{equation}
2=[\CC:K]=2[\RR:K],
\end{equation}
which implies $K=\RR$.
\end{proof}

\subsection{Proof of the Artin--Schreier Theorem}

\begin{recall}
The theorem statement: Let $L/F$ be a field extension.
If $L$ is an algebraically closed field with $1<[L:F]<\infty$, then
$L=F(\I)$ where $\I=\sqrt{-1}$ and $\Char(F)=0$.
\end{recall}

\begin{claim}
$L/K$ is Galois.
\end{claim}

\begin{proof}
We see that $L/K$ is normal (since $L$ is algebraically closed).

To prove $L/K$ is separable, either $\Char(F)=0$ or $\Char(F)=p>0$.

\textsc{Case 1}: If $\Char(F)=0$, then it's perfect and we are good.

\textsc{Case 2}: If $\Char(F)=p>0$, then we will show $F^{p}=F$. If
$a\in F\setminus F^{p}$, $x^{p^{m}}-a$ is irreducible for all
$m\in\NN$. This means $F$ has extension of degree $p^{m}$ for all
$m\in\NN$ which contradicts $[L:F]<\infty$. Hence $L$ mst be separable.
\end{proof}

\begin{claim}
There is no odd prime number $p$ such that $p\divides[L:F]$.
\end{claim}

\begin{proof}
Assume for contradiction there exists some odd prime $p$ such that
$p\divides[L:F]$. Then since $|\Gal(L/F)|=[L:F]$, there is a subgroup
of $\Gal(L/F)$ of order $p$ (hence cyclic). Then there exists a
intermediate subextension $L/K/F$ such that $\Gal(L/K)$ is cyclic of
order $p$. We see
\begin{equation}
\Char(F)=\Char(K)\neq p.
\end{equation}
If not, by Proposition~\ref{prop:11.3} we have $L=K(\alpha)$ where
$\alpha$ is a root of $x^{p}-x-a\in K[x]$. Then
$\{1,\alpha,\dots,\alpha^{p-1}\}$ is a basis of $L/K$. For any $b\in
L$, we have
\begin{equation}
b=b_{0}+b_{1}\alpha+\dots+b_{p-1}\alpha^{p-1}.
\end{equation}
Then we can write
\begin{equation}
b^{p}-b=\sum^{p-1}_{j=0}(b^{p}_{j}(\alpha+a)^{j}-b_{j}\alpha^{j}).
\end{equation}
Look at the coefficient of $\alpha^{p-1}$: $b^{p}_{p-1}-b_{p-1}$.
Since $L$ is algebraically closed, we can pick $b$ such that
$b^{p}-b=a\alpha^{p-1}$. Then
\begin{equation}
b^{p}_{p-1}-b_{p-1}-a=0.
\end{equation}
But this is a contradiction, hence $\Char(F)\neq p$.
\end{proof}

\begin{claim}
$[L:F]=2$.
\end{claim}

